% ============================================================
% Mục 6. Mô hình toán học  (Đề tài 9)
% Đây là đoạn nội dung (không có preamble): dùng % ============================================================
% Mục 6. Mô hình toán học  (Đề tài 9)
% Đây là đoạn nội dung (không có preamble): dùng % ============================================================
% Mục 6. Mô hình toán học  (Đề tài 9)
% Đây là đoạn nội dung (không có preamble): dùng % ============================================================
% Mục 6. Mô hình toán học  (Đề tài 9)
% Đây là đoạn nội dung (không có preamble): dùng \input{muc6} hoặc dán vào báo cáo.
% Cần: amsmath, booktabs. Số liệu khớp với src/model.py.
% ============================================================
\section{Mô hình toán học}\label{sec:mohinh}

\subsection{Bối cảnh và đại lượng khảo sát}
Xét một doanh nghiệp bán một loại sản phẩm với giá bán $p$. Giá bán là đại lượng
doanh nghiệp \emph{dự kiến} khi lập kế hoạch, nhưng giá thực tế thu được thường lệch
so với dự kiến (do chiết khấu, khuyến mãi, làm tròn giá, hoặc biến động thị trường).
Bài toán đặt ra: \emph{sai số nhỏ của giá làm doanh thu và lợi nhuận thay đổi bao
nhiêu, và có thể dự báo nhanh sự thay đổi đó bằng vi phân đến mức độ chính xác nào?}

Các đại lượng của mô hình được cho trong Bảng~\ref{tab:var}.

\begin{table}[h]
\centering
\caption{Các đại lượng trong mô hình và đơn vị.}\label{tab:var}
\begin{tabular}{@{}clll@{}}
\toprule
Ký hiệu & Ý nghĩa & Đơn vị & Vai trò \\
\midrule
$p$ & giá bán & nghìn đồng/sản phẩm & biến đầu vào \\
$q(p)$ & lượng cầu (số sản phẩm bán được) & sản phẩm/tuần & đầu ra \\
$R(p)$ & doanh thu & nghìn đồng/tuần & đầu ra \\
$P(p)$ & lợi nhuận & nghìn đồng/tuần & đầu ra \\
$c$ & chi phí biên (giá vốn) mỗi sản phẩm & nghìn đồng/sản phẩm & tham số \\
\bottomrule
\end{tabular}
\end{table}

Doanh thu và lợi nhuận được xác định từ lượng cầu:
\begin{equation}
R(p) = p\,q(p), \qquad P(p) = R(p) - c\,q(p) = (p-c)\,q(p).
\end{equation}
Trong đó $c$ được giả thiết không đổi và không phụ thuộc sản lượng.

\subsection{Hai mô hình cầu}
Để xem kết luận phụ thuộc thế nào vào dạng hàm cầu, nhóm khảo sát hai mô hình với
tính chất khác nhau.

\paragraph{Mô hình (L): cầu tuyến tính.}
\begin{equation}
q(p) = a - b\,p, \qquad 0 < p < \frac{a}{b},
\end{equation}
với $a>0$ là lượng cầu tối đa khi giá tiến về $0$ và $b>0$ là số sản phẩm cầu giảm đi
khi giá tăng thêm $1$ nghìn đồng. Điều kiện $p<a/b$ bảo đảm $q>0$. Khi đó
\begin{align}
R(p) &= ap - bp^2, \\
P(p) &= (p-c)(a-bp) = -bp^2 + (a+bc)\,p - ac.
\end{align}

\paragraph{Mô hình (E): cầu co giãn không đổi.}
\begin{equation}
q(p) = A\,p^{-k}, \qquad p>0,\ k>1,
\end{equation}
với $A>0$ là hằng số quy mô và $k$ là độ co giãn của cầu theo giá (giá tăng $1\%$ thì
cầu giảm khoảng $k\%$). Điều kiện $k>1$ bảo đảm lợi nhuận có giá tối ưu hữu hạn. Khi đó
\begin{align}
R(p) &= A\,p^{1-k}, \\
P(p) &= A\,p^{1-k} - cA\,p^{-k}.
\end{align}

\subsection{Giá tối ưu}
Giá tối ưu $p^*$ là điểm làm lợi nhuận lớn nhất, thỏa $P'(p^*)=0$ (cách tính ở
Mục~\ref{sec:phuongphap}):
\begin{equation}
\text{(L)}\quad p^* = \frac{a+bc}{2b}, \qquad\qquad
\text{(E)}\quad p^* = \frac{ck}{k-1}.
\end{equation}
Điểm $p^*$ có vai trò đặc biệt trong bài toán độ nhạy: tại đó đạo hàm bậc nhất của $P$
bằng $0$, nên vi phân bậc nhất không còn phản ánh đúng sự thay đổi của lợi nhuận
(xem Mục~\ref{sec:thucnghiem}).

\subsection{Bộ tham số của ví dụ}
Nhóm sử dụng bộ tham số giả định, do nhóm tự xây dựng (không lấy từ giáo trình), với
chi phí biên $c=30$ nghìn đồng/sản phẩm:

\begin{table}[h]
\centering
\caption{Tham số và điểm khảo sát.}\label{tab:par}
\begin{tabular}{@{}lccccc@{}}
\toprule
Mô hình & Tham số & Giá hiện tại $p_0$ & $q(p_0)$ & $R(p_0)$ & $P(p_0)$ \\
\midrule
(L) & $a=1200,\ b=10$ & $50$ & $700$ & $35\,000$ & $14\,000$ \\
(E) & $A=10^6,\ k=2$ & $40$ & $625$ & $25\,000$ & $6\,250$ \\
\bottomrule
\end{tabular}
\end{table}

Giá tối ưu tương ứng là $p^*=75$ (mô hình L, $P(p^*)=20\,250$) và $p^*=60$ (mô hình E,
$P(p^*)\approx 8\,333{,}3$), đơn vị lợi nhuận là nghìn đồng/tuần. Hai bộ số cho hai
tình huống khác nhau: ở (L) giá hiện tại thấp hơn giá tối ưu và cầu giảm đều theo giá,
ở (E) cầu giảm rất nhanh khi giá thấp và chậm dần khi giá cao.

\subsection{Sai số đầu vào và câu hỏi nghiên cứu}
Giả sử giá thực tế là $p_0+\Delta p$ thay vì $p_0$. Sai số tương đối của đầu vào là
$\Delta p/p_0$ (ví dụ $\pm 1\%,\ \pm 5\%,\ \pm 20\%$). Với mỗi đại lượng đầu ra
$y=f(p)\in\{q,R,P\}$, thay đổi \emph{thực} là
\begin{equation}
\Delta y = f(p_0+\Delta p) - f(p_0).
\end{equation}
Nhóm trả lời bốn câu hỏi, tương ứng các yêu cầu của đề tài:
\begin{enumerate}
  \item Độ nhạy của $q$, $R$, $P$ theo giá tại $p_0$ lớn hay nhỏ, và đại lượng nào nhạy nhất?
  \item Vi phân $dy=f'(p_0)\Delta p$ dự báo $\Delta y$ chính xác đến mức nào khi
        $|\Delta p/p_0|$ tăng từ $1\%$ lên $20\%$?
  \item Sai số của xấp xỉ $|\Delta y - dy|$ phụ thuộc vào $\Delta p$ ra sao?
  \item Khi giá hiện tại ở gần giá tối ưu, kết luận thay đổi thế nào?
\end{enumerate}

\subsection{Giả thiết và hạn chế của mô hình}
\begin{itemize}
  \item Chỉ giá $p$ bị sai số; các tham số $a,b,A,k,c$ được coi là chính xác.
  \item Lượng cầu phản ứng tức thời và tất định theo giá; không xét yếu tố mùa vụ,
        đối thủ cạnh tranh hay giới hạn sản xuất.
  \item Chi phí biên $c$ không đổi.
  \item Bộ tham số là giả định minh họa, không được ước lượng từ dữ liệu thị trường.
        Kết luận về \emph{phương pháp} (cách sai số lan truyền) có tính tổng quát,
        còn các con số cụ thể chỉ có ý nghĩa với bộ tham số này.
\end{itemize}
 hoặc dán vào báo cáo.
% Cần: amsmath, booktabs. Số liệu khớp với src/model.py.
% ============================================================
\section{Mô hình toán học}\label{sec:mohinh}

\subsection{Bối cảnh và đại lượng khảo sát}
Xét một doanh nghiệp bán một loại sản phẩm với giá bán $p$. Giá bán là đại lượng
doanh nghiệp \emph{dự kiến} khi lập kế hoạch, nhưng giá thực tế thu được thường lệch
so với dự kiến (do chiết khấu, khuyến mãi, làm tròn giá, hoặc biến động thị trường).
Bài toán đặt ra: \emph{sai số nhỏ của giá làm doanh thu và lợi nhuận thay đổi bao
nhiêu, và có thể dự báo nhanh sự thay đổi đó bằng vi phân đến mức độ chính xác nào?}

Các đại lượng của mô hình được cho trong Bảng~\ref{tab:var}.

\begin{table}[h]
\centering
\caption{Các đại lượng trong mô hình và đơn vị.}\label{tab:var}
\begin{tabular}{@{}clll@{}}
\toprule
Ký hiệu & Ý nghĩa & Đơn vị & Vai trò \\
\midrule
$p$ & giá bán & nghìn đồng/sản phẩm & biến đầu vào \\
$q(p)$ & lượng cầu (số sản phẩm bán được) & sản phẩm/tuần & đầu ra \\
$R(p)$ & doanh thu & nghìn đồng/tuần & đầu ra \\
$P(p)$ & lợi nhuận & nghìn đồng/tuần & đầu ra \\
$c$ & chi phí biên (giá vốn) mỗi sản phẩm & nghìn đồng/sản phẩm & tham số \\
\bottomrule
\end{tabular}
\end{table}

Doanh thu và lợi nhuận được xác định từ lượng cầu:
\begin{equation}
R(p) = p\,q(p), \qquad P(p) = R(p) - c\,q(p) = (p-c)\,q(p).
\end{equation}
Trong đó $c$ được giả thiết không đổi và không phụ thuộc sản lượng.

\subsection{Hai mô hình cầu}
Để xem kết luận phụ thuộc thế nào vào dạng hàm cầu, nhóm khảo sát hai mô hình với
tính chất khác nhau.

\paragraph{Mô hình (L): cầu tuyến tính.}
\begin{equation}
q(p) = a - b\,p, \qquad 0 < p < \frac{a}{b},
\end{equation}
với $a>0$ là lượng cầu tối đa khi giá tiến về $0$ và $b>0$ là số sản phẩm cầu giảm đi
khi giá tăng thêm $1$ nghìn đồng. Điều kiện $p<a/b$ bảo đảm $q>0$. Khi đó
\begin{align}
R(p) &= ap - bp^2, \\
P(p) &= (p-c)(a-bp) = -bp^2 + (a+bc)\,p - ac.
\end{align}

\paragraph{Mô hình (E): cầu co giãn không đổi.}
\begin{equation}
q(p) = A\,p^{-k}, \qquad p>0,\ k>1,
\end{equation}
với $A>0$ là hằng số quy mô và $k$ là độ co giãn của cầu theo giá (giá tăng $1\%$ thì
cầu giảm khoảng $k\%$). Điều kiện $k>1$ bảo đảm lợi nhuận có giá tối ưu hữu hạn. Khi đó
\begin{align}
R(p) &= A\,p^{1-k}, \\
P(p) &= A\,p^{1-k} - cA\,p^{-k}.
\end{align}

\subsection{Giá tối ưu}
Giá tối ưu $p^*$ là điểm làm lợi nhuận lớn nhất, thỏa $P'(p^*)=0$ (cách tính ở
Mục~\ref{sec:phuongphap}):
\begin{equation}
\text{(L)}\quad p^* = \frac{a+bc}{2b}, \qquad\qquad
\text{(E)}\quad p^* = \frac{ck}{k-1}.
\end{equation}
Điểm $p^*$ có vai trò đặc biệt trong bài toán độ nhạy: tại đó đạo hàm bậc nhất của $P$
bằng $0$, nên vi phân bậc nhất không còn phản ánh đúng sự thay đổi của lợi nhuận
(xem Mục~\ref{sec:thucnghiem}).

\subsection{Bộ tham số của ví dụ}
Nhóm sử dụng bộ tham số giả định, do nhóm tự xây dựng (không lấy từ giáo trình), với
chi phí biên $c=30$ nghìn đồng/sản phẩm:

\begin{table}[h]
\centering
\caption{Tham số và điểm khảo sát.}\label{tab:par}
\begin{tabular}{@{}lccccc@{}}
\toprule
Mô hình & Tham số & Giá hiện tại $p_0$ & $q(p_0)$ & $R(p_0)$ & $P(p_0)$ \\
\midrule
(L) & $a=1200,\ b=10$ & $50$ & $700$ & $35\,000$ & $14\,000$ \\
(E) & $A=10^6,\ k=2$ & $40$ & $625$ & $25\,000$ & $6\,250$ \\
\bottomrule
\end{tabular}
\end{table}

Giá tối ưu tương ứng là $p^*=75$ (mô hình L, $P(p^*)=20\,250$) và $p^*=60$ (mô hình E,
$P(p^*)\approx 8\,333{,}3$), đơn vị lợi nhuận là nghìn đồng/tuần. Hai bộ số cho hai
tình huống khác nhau: ở (L) giá hiện tại thấp hơn giá tối ưu và cầu giảm đều theo giá,
ở (E) cầu giảm rất nhanh khi giá thấp và chậm dần khi giá cao.

\subsection{Sai số đầu vào và câu hỏi nghiên cứu}
Giả sử giá thực tế là $p_0+\Delta p$ thay vì $p_0$. Sai số tương đối của đầu vào là
$\Delta p/p_0$ (ví dụ $\pm 1\%,\ \pm 5\%,\ \pm 20\%$). Với mỗi đại lượng đầu ra
$y=f(p)\in\{q,R,P\}$, thay đổi \emph{thực} là
\begin{equation}
\Delta y = f(p_0+\Delta p) - f(p_0).
\end{equation}
Nhóm trả lời bốn câu hỏi, tương ứng các yêu cầu của đề tài:
\begin{enumerate}
  \item Độ nhạy của $q$, $R$, $P$ theo giá tại $p_0$ lớn hay nhỏ, và đại lượng nào nhạy nhất?
  \item Vi phân $dy=f'(p_0)\Delta p$ dự báo $\Delta y$ chính xác đến mức nào khi
        $|\Delta p/p_0|$ tăng từ $1\%$ lên $20\%$?
  \item Sai số của xấp xỉ $|\Delta y - dy|$ phụ thuộc vào $\Delta p$ ra sao?
  \item Khi giá hiện tại ở gần giá tối ưu, kết luận thay đổi thế nào?
\end{enumerate}

\subsection{Giả thiết và hạn chế của mô hình}
\begin{itemize}
  \item Chỉ giá $p$ bị sai số; các tham số $a,b,A,k,c$ được coi là chính xác.
  \item Lượng cầu phản ứng tức thời và tất định theo giá; không xét yếu tố mùa vụ,
        đối thủ cạnh tranh hay giới hạn sản xuất.
  \item Chi phí biên $c$ không đổi.
  \item Bộ tham số là giả định minh họa, không được ước lượng từ dữ liệu thị trường.
        Kết luận về \emph{phương pháp} (cách sai số lan truyền) có tính tổng quát,
        còn các con số cụ thể chỉ có ý nghĩa với bộ tham số này.
\end{itemize}
 hoặc dán vào báo cáo.
% Cần: amsmath, booktabs. Số liệu khớp với src/model.py.
% ============================================================
\section{Mô hình toán học}\label{sec:mohinh}

\subsection{Bối cảnh và đại lượng khảo sát}
Xét một doanh nghiệp bán một loại sản phẩm với giá bán $p$. Giá bán là đại lượng
doanh nghiệp \emph{dự kiến} khi lập kế hoạch, nhưng giá thực tế thu được thường lệch
so với dự kiến (do chiết khấu, khuyến mãi, làm tròn giá, hoặc biến động thị trường).
Bài toán đặt ra: \emph{sai số nhỏ của giá làm doanh thu và lợi nhuận thay đổi bao
nhiêu, và có thể dự báo nhanh sự thay đổi đó bằng vi phân đến mức độ chính xác nào?}

Các đại lượng của mô hình được cho trong Bảng~\ref{tab:var}.

\begin{table}[h]
\centering
\caption{Các đại lượng trong mô hình và đơn vị.}\label{tab:var}
\begin{tabular}{@{}clll@{}}
\toprule
Ký hiệu & Ý nghĩa & Đơn vị & Vai trò \\
\midrule
$p$ & giá bán & nghìn đồng/sản phẩm & biến đầu vào \\
$q(p)$ & lượng cầu (số sản phẩm bán được) & sản phẩm/tuần & đầu ra \\
$R(p)$ & doanh thu & nghìn đồng/tuần & đầu ra \\
$P(p)$ & lợi nhuận & nghìn đồng/tuần & đầu ra \\
$c$ & chi phí biên (giá vốn) mỗi sản phẩm & nghìn đồng/sản phẩm & tham số \\
\bottomrule
\end{tabular}
\end{table}

Doanh thu và lợi nhuận được xác định từ lượng cầu:
\begin{equation}
R(p) = p\,q(p), \qquad P(p) = R(p) - c\,q(p) = (p-c)\,q(p).
\end{equation}
Trong đó $c$ được giả thiết không đổi và không phụ thuộc sản lượng.

\subsection{Hai mô hình cầu}
Để xem kết luận phụ thuộc thế nào vào dạng hàm cầu, nhóm khảo sát hai mô hình với
tính chất khác nhau.

\paragraph{Mô hình (L): cầu tuyến tính.}
\begin{equation}
q(p) = a - b\,p, \qquad 0 < p < \frac{a}{b},
\end{equation}
với $a>0$ là lượng cầu tối đa khi giá tiến về $0$ và $b>0$ là số sản phẩm cầu giảm đi
khi giá tăng thêm $1$ nghìn đồng. Điều kiện $p<a/b$ bảo đảm $q>0$. Khi đó
\begin{align}
R(p) &= ap - bp^2, \\
P(p) &= (p-c)(a-bp) = -bp^2 + (a+bc)\,p - ac.
\end{align}

\paragraph{Mô hình (E): cầu co giãn không đổi.}
\begin{equation}
q(p) = A\,p^{-k}, \qquad p>0,\ k>1,
\end{equation}
với $A>0$ là hằng số quy mô và $k$ là độ co giãn của cầu theo giá (giá tăng $1\%$ thì
cầu giảm khoảng $k\%$). Điều kiện $k>1$ bảo đảm lợi nhuận có giá tối ưu hữu hạn. Khi đó
\begin{align}
R(p) &= A\,p^{1-k}, \\
P(p) &= A\,p^{1-k} - cA\,p^{-k}.
\end{align}

\subsection{Giá tối ưu}
Giá tối ưu $p^*$ là điểm làm lợi nhuận lớn nhất, thỏa $P'(p^*)=0$ (cách tính ở
Mục~\ref{sec:phuongphap}):
\begin{equation}
\text{(L)}\quad p^* = \frac{a+bc}{2b}, \qquad\qquad
\text{(E)}\quad p^* = \frac{ck}{k-1}.
\end{equation}
Điểm $p^*$ có vai trò đặc biệt trong bài toán độ nhạy: tại đó đạo hàm bậc nhất của $P$
bằng $0$, nên vi phân bậc nhất không còn phản ánh đúng sự thay đổi của lợi nhuận
(xem Mục~\ref{sec:thucnghiem}).

\subsection{Bộ tham số của ví dụ}
Nhóm sử dụng bộ tham số giả định, do nhóm tự xây dựng (không lấy từ giáo trình), với
chi phí biên $c=30$ nghìn đồng/sản phẩm:

\begin{table}[h]
\centering
\caption{Tham số và điểm khảo sát.}\label{tab:par}
\begin{tabular}{@{}lccccc@{}}
\toprule
Mô hình & Tham số & Giá hiện tại $p_0$ & $q(p_0)$ & $R(p_0)$ & $P(p_0)$ \\
\midrule
(L) & $a=1200,\ b=10$ & $50$ & $700$ & $35\,000$ & $14\,000$ \\
(E) & $A=10^6,\ k=2$ & $40$ & $625$ & $25\,000$ & $6\,250$ \\
\bottomrule
\end{tabular}
\end{table}

Giá tối ưu tương ứng là $p^*=75$ (mô hình L, $P(p^*)=20\,250$) và $p^*=60$ (mô hình E,
$P(p^*)\approx 8\,333{,}3$), đơn vị lợi nhuận là nghìn đồng/tuần. Hai bộ số cho hai
tình huống khác nhau: ở (L) giá hiện tại thấp hơn giá tối ưu và cầu giảm đều theo giá,
ở (E) cầu giảm rất nhanh khi giá thấp và chậm dần khi giá cao.

\subsection{Sai số đầu vào và câu hỏi nghiên cứu}
Giả sử giá thực tế là $p_0+\Delta p$ thay vì $p_0$. Sai số tương đối của đầu vào là
$\Delta p/p_0$ (ví dụ $\pm 1\%,\ \pm 5\%,\ \pm 20\%$). Với mỗi đại lượng đầu ra
$y=f(p)\in\{q,R,P\}$, thay đổi \emph{thực} là
\begin{equation}
\Delta y = f(p_0+\Delta p) - f(p_0).
\end{equation}
Nhóm trả lời bốn câu hỏi, tương ứng các yêu cầu của đề tài:
\begin{enumerate}
  \item Độ nhạy của $q$, $R$, $P$ theo giá tại $p_0$ lớn hay nhỏ, và đại lượng nào nhạy nhất?
  \item Vi phân $dy=f'(p_0)\Delta p$ dự báo $\Delta y$ chính xác đến mức nào khi
        $|\Delta p/p_0|$ tăng từ $1\%$ lên $20\%$?
  \item Sai số của xấp xỉ $|\Delta y - dy|$ phụ thuộc vào $\Delta p$ ra sao?
  \item Khi giá hiện tại ở gần giá tối ưu, kết luận thay đổi thế nào?
\end{enumerate}

\subsection{Giả thiết và hạn chế của mô hình}
\begin{itemize}
  \item Chỉ giá $p$ bị sai số; các tham số $a,b,A,k,c$ được coi là chính xác.
  \item Lượng cầu phản ứng tức thời và tất định theo giá; không xét yếu tố mùa vụ,
        đối thủ cạnh tranh hay giới hạn sản xuất.
  \item Chi phí biên $c$ không đổi.
  \item Bộ tham số là giả định minh họa, không được ước lượng từ dữ liệu thị trường.
        Kết luận về \emph{phương pháp} (cách sai số lan truyền) có tính tổng quát,
        còn các con số cụ thể chỉ có ý nghĩa với bộ tham số này.
\end{itemize}
 hoặc dán vào báo cáo.
% Cần: amsmath, booktabs. Số liệu khớp với src/model.py.
% ============================================================
\section{Mô hình toán học}\label{sec:mohinh}

\subsection{Bối cảnh và đại lượng khảo sát}
Xét một doanh nghiệp bán một loại sản phẩm với giá bán $p$. Giá bán là đại lượng
doanh nghiệp \emph{dự kiến} khi lập kế hoạch, nhưng giá thực tế thu được thường lệch
so với dự kiến (do chiết khấu, khuyến mãi, làm tròn giá, hoặc biến động thị trường).
Bài toán đặt ra: \emph{sai số nhỏ của giá làm doanh thu và lợi nhuận thay đổi bao
nhiêu, và có thể dự báo nhanh sự thay đổi đó bằng vi phân đến mức độ chính xác nào?}

Các đại lượng của mô hình được cho trong Bảng~\ref{tab:var}.

\begin{table}[h]
\centering
\caption{Các đại lượng trong mô hình và đơn vị.}\label{tab:var}
\begin{tabular}{@{}clll@{}}
\toprule
Ký hiệu & Ý nghĩa & Đơn vị & Vai trò \\
\midrule
$p$ & giá bán & nghìn đồng/sản phẩm & biến đầu vào \\
$q(p)$ & lượng cầu (số sản phẩm bán được) & sản phẩm/tuần & đầu ra \\
$R(p)$ & doanh thu & nghìn đồng/tuần & đầu ra \\
$P(p)$ & lợi nhuận & nghìn đồng/tuần & đầu ra \\
$c$ & chi phí biên (giá vốn) mỗi sản phẩm & nghìn đồng/sản phẩm & tham số \\
\bottomrule
\end{tabular}
\end{table}

Doanh thu và lợi nhuận được xác định từ lượng cầu:
\begin{equation}
R(p) = p\,q(p), \qquad P(p) = R(p) - c\,q(p) = (p-c)\,q(p).
\end{equation}
Trong đó $c$ được giả thiết không đổi và không phụ thuộc sản lượng.

\subsection{Hai mô hình cầu}
Để xem kết luận phụ thuộc thế nào vào dạng hàm cầu, nhóm khảo sát hai mô hình với
tính chất khác nhau.

\paragraph{Mô hình (L): cầu tuyến tính.}
\begin{equation}
q(p) = a - b\,p, \qquad 0 < p < \frac{a}{b},
\end{equation}
với $a>0$ là lượng cầu tối đa khi giá tiến về $0$ và $b>0$ là số sản phẩm cầu giảm đi
khi giá tăng thêm $1$ nghìn đồng. Điều kiện $p<a/b$ bảo đảm $q>0$. Khi đó
\begin{align}
R(p) &= ap - bp^2, \\
P(p) &= (p-c)(a-bp) = -bp^2 + (a+bc)\,p - ac.
\end{align}

\paragraph{Mô hình (E): cầu co giãn không đổi.}
\begin{equation}
q(p) = A\,p^{-k}, \qquad p>0,\ k>1,
\end{equation}
với $A>0$ là hằng số quy mô và $k$ là độ co giãn của cầu theo giá (giá tăng $1\%$ thì
cầu giảm khoảng $k\%$). Điều kiện $k>1$ bảo đảm lợi nhuận có giá tối ưu hữu hạn. Khi đó
\begin{align}
R(p) &= A\,p^{1-k}, \\
P(p) &= A\,p^{1-k} - cA\,p^{-k}.
\end{align}

\subsection{Giá tối ưu}
Giá tối ưu $p^*$ là điểm làm lợi nhuận lớn nhất, thỏa $P'(p^*)=0$ (cách tính ở
Mục~\ref{sec:phuongphap}):
\begin{equation}
\text{(L)}\quad p^* = \frac{a+bc}{2b}, \qquad\qquad
\text{(E)}\quad p^* = \frac{ck}{k-1}.
\end{equation}
Điểm $p^*$ có vai trò đặc biệt trong bài toán độ nhạy: tại đó đạo hàm bậc nhất của $P$
bằng $0$, nên vi phân bậc nhất không còn phản ánh đúng sự thay đổi của lợi nhuận
(xem Mục~\ref{sec:thucnghiem}).

\subsection{Bộ tham số của ví dụ}
Nhóm sử dụng bộ tham số giả định, do nhóm tự xây dựng (không lấy từ giáo trình), với
chi phí biên $c=30$ nghìn đồng/sản phẩm:

\begin{table}[h]
\centering
\caption{Tham số và điểm khảo sát.}\label{tab:par}
\begin{tabular}{@{}lccccc@{}}
\toprule
Mô hình & Tham số & Giá hiện tại $p_0$ & $q(p_0)$ & $R(p_0)$ & $P(p_0)$ \\
\midrule
(L) & $a=1200,\ b=10$ & $50$ & $700$ & $35\,000$ & $14\,000$ \\
(E) & $A=10^6,\ k=2$ & $40$ & $625$ & $25\,000$ & $6\,250$ \\
\bottomrule
\end{tabular}
\end{table}

Giá tối ưu tương ứng là $p^*=75$ (mô hình L, $P(p^*)=20\,250$) và $p^*=60$ (mô hình E,
$P(p^*)\approx 8\,333{,}3$), đơn vị lợi nhuận là nghìn đồng/tuần. Hai bộ số cho hai
tình huống khác nhau: ở (L) giá hiện tại thấp hơn giá tối ưu và cầu giảm đều theo giá,
ở (E) cầu giảm rất nhanh khi giá thấp và chậm dần khi giá cao.

\subsection{Sai số đầu vào và câu hỏi nghiên cứu}
Giả sử giá thực tế là $p_0+\Delta p$ thay vì $p_0$. Sai số tương đối của đầu vào là
$\Delta p/p_0$ (ví dụ $\pm 1\%,\ \pm 5\%,\ \pm 20\%$). Với mỗi đại lượng đầu ra
$y=f(p)\in\{q,R,P\}$, thay đổi \emph{thực} là
\begin{equation}
\Delta y = f(p_0+\Delta p) - f(p_0).
\end{equation}
Nhóm trả lời bốn câu hỏi, tương ứng các yêu cầu của đề tài:
\begin{enumerate}
  \item Độ nhạy của $q$, $R$, $P$ theo giá tại $p_0$ lớn hay nhỏ, và đại lượng nào nhạy nhất?
  \item Vi phân $dy=f'(p_0)\Delta p$ dự báo $\Delta y$ chính xác đến mức nào khi
        $|\Delta p/p_0|$ tăng từ $1\%$ lên $20\%$?
  \item Sai số của xấp xỉ $|\Delta y - dy|$ phụ thuộc vào $\Delta p$ ra sao?
  \item Khi giá hiện tại ở gần giá tối ưu, kết luận thay đổi thế nào?
\end{enumerate}

\subsection{Giả thiết và hạn chế của mô hình}
\begin{itemize}
  \item Chỉ giá $p$ bị sai số; các tham số $a,b,A,k,c$ được coi là chính xác.
  \item Lượng cầu phản ứng tức thời và tất định theo giá; không xét yếu tố mùa vụ,
        đối thủ cạnh tranh hay giới hạn sản xuất.
  \item Chi phí biên $c$ không đổi.
  \item Bộ tham số là giả định minh họa, không được ước lượng từ dữ liệu thị trường.
        Kết luận về \emph{phương pháp} (cách sai số lan truyền) có tính tổng quát,
        còn các con số cụ thể chỉ có ý nghĩa với bộ tham số này.
\end{itemize}
